% engines: lualatex, pdflatex
% expect-warnings: 0
% expect-errors: 0
\documentclass{article}
\usepackage[margin=1cm]{geometry}
\usepackage{../../latex/pgf-go}
\newcommand\case[1]{\par\medskip\noindent\textbf{#1}\par\nopagebreak}
\pagestyle{empty}
\begin{document}

% The layers of the goban have no scopes: what a mark sets in the PDF must not reach the marks
% drawn after it. And a shape key forgets what the previous shape key set.

\case{dim (B1) and then opaque filled marks (C1, D1, C2) and a stone mark (D2)}
\begin{goban}[5x2]
	\stones[b]{D2}
	\marks[dim]{B1}
	\marks[filled square]{C1}\marks[filled circle]{D1}\marks[circle]{C2}\marks[w, triangle]{D2}
\end{goban}

\case{Translucent grid: t-circle (B1), then opaque circle and cross (C1, D1)}
\begin{goban}[5x1, grid = {line = {opacity = 0.3}}]
	\marks[t-circle]{B1}
	\marks[circle]{C1}\marks[cross]{D1}
\end{goban}

\case{t-circle, circle (B1) and hatched lines, circle (D1) are plain circles}
\begin{goban}[5x1]
	\marks[t-circle, circle]{B1}
	\marks[circle]{C1}
	\marks[hatched lines, circle]{D1}
\end{goban}

\end{document}
